\clearpage\begingroup
\hypertarget{native-333b}{}\HistoricalHead{328}{L}{PRAKTISCHE UND SPHÄRISCHE ASTRONOMIE. BRIEFWECHSEL.}
\begin{center}\textsc{Gauss} an \textsc{Bessel}, 7. December 1820.\end{center}
\noindent\hspace*{\parindent}. . . . . . Bei der Reduction der Beobachtungen zu Polardistanzen hingegen nehme ich den Ort des Pols nur von einer Umlegung zur nächsten als
\clearpage\SourceLeaf{334}\HistoricalHead{329}{R}{BEOBACHTUNGEN AM REICHENBACHSCHEN MERIDIANKREISE.}\HistoricalFoot{XI1.}{42}
\noindent constant an und bestimme diesen aus dem Ensemble \textls[120]{aller} Circumpolarsterne (etwa vom Zenith bis zum Pol), die während der Zeit in beiden Culminationen beobachtet sind. Wenn dies Verfahren immer strenge befolgt wird, so ist wenigstens am Ende kein constanter Fehler zu befürchten. Meine Beobachtungen vom Febr[uar] bis Sept[ember] zerfallen so in 7 Perioden, und ich hatte bis dahin erst für die 4 ersten Perioden den Ort des Pols definitiv bestimmt. Zur Reduction der übrigen 3 konnte ich lange nach meiner Zurückkunft, wo eine Menge zum Theil verdriesslicher Geschäfte auf mich warten, keine Zeit finden. Ich habe daher, um, was noch von jener übrig war, abzuthun, auch noch die beiden letzten Nächte zu Hülfe genommen und werde Ihnen sogleich meine Endresultate anzeigen. . . . . . .
\begin{center}\begingroup\setlength{\tabcolsep}{5pt}
\begin{tabular}{@{}lrrrc@{}}
\multicolumn{4}{c}{Ort des Pols in den einzelnen Perioden:}&Gewicht:\\
I.&\(321^\circ\)&\(29'\)&\(31\overset{\prime\prime}{,}92\)&\(43{,}9\)\\
II.&38&25&\(53{,}58\)&\(64{,}0\)\\
III.&321&29&\(30{,}78\)&\(90{,}8\)\\
IV.&38&25&\(53{,}68\)&\(25{,}1\)\\
V.&321&29&\(31{,}37\)&\(21{,}5\)\\
VI.&38&25&\(54{,}53\)&\(10{,}7\)\\
VII.&321&29&\(32{,}74\)&\(8{,}7\)
\end{tabular}\par\endgroup\end{center}

\noindent\hspace*{\parindent}Die einzelnen Bestimmungen stimmen immer sehr gut und zeigen gar keine Spur eines Unterschiedes nach den Declinationen, als Beweis der Vortrefflichkeit Ihrer Refractionstafel; ja selbst noch südlichere Sterne, wie \(\alpha\) Cygni, \(\alpha\) Aurigae, stimmen immer gut genug, ich habe sie aber nie zugezogen. Als Beispiel die fünfte Periode:
\begin{center}\begingroup\setlength{\tabcolsep}{5pt}
\begin{tabular}{@{}lrlr@{}}
\(\delta\) Ursae min.&\(321^\circ29'30\overset{\prime\prime}{,}64\)&Gewicht:&\(2{,}0\)\\
\(\varepsilon\) Ursae min.&\(31{,}41\)&&\(4{,}0\)\\
1 Drac[onis] Hev.&\(31{,}70\)&&\(2{,}0\)\\
\(\zeta\) Dracon[is]&\(31{,}13\)&&\(2{,}7\)\\
\(\mathrm{o}'\) Ursae mai.&\(31{,}57\)&&\(2{,}0\)\\
\(d\) Camelop[ardalis]&\(31{,}41\)&&\(4{,}8\)\\
\(\delta\) Aurigae&\(31{,}44\)&&\(2{,}0\)\\
\(\theta\) Ursae mai.&\(31{,}65\)&&\(2{,}0\)
\end{tabular}\par\endgroup\end{center}

\noindent\hspace*{\parindent}Die Reductionen der Sterne sind nun so gemacht, dass ich, wenn der
\clearpage\SourceLeaf{335}\HistoricalHead{330}{L}{PRAKTISCHE UND SPHÄRISCHE ASTRONOMIE. BRIEFWECHSEL.}
\noindent Stern in einer Periode \(a\)-mal beob[achtet] ist und das Gewicht der Best[immung] des Pols \(=b\) war, das Gewicht der Declination \(\dfrac{ab}{a+b}\) setze. Z. B. [Piazzi] XX, 222 ist in der fünften Periode einmal, in der sechsten dreimal, in der siebenten dreimal beobachtet, daraus
\begin{center}\begin{tabular}{@{}lrlr@{}}
V.&\(59^\circ48'55\overset{\prime\prime}{,}73\)&Gewicht&\(1{,}0\)\\
VI.&\(56{,}16\)&&\(2{,}3\)\\
VII.&\(56{,}60\)&&\(2{,}2\)\\\hline
&\(59^\circ48'56\overset{\prime\prime}{,}26\)&&\(5{,}5\)
\end{tabular}\end{center}

\noindent\hspace*{\parindent}Meistens stimmen die Resultate[\begingroup\renewcommand{\thefootnote}{*)}\footnotemark\renewcommand{\thefootnote}{[*)}\footnotetext{Für die von \textsc{Schumacher} in Lauenburg beobachteten Sterne; vergl. den Brief an \textsc{Bessel} vom 17. August 1820, S. 327.]}\endgroup] aus verschiedenen Perioden recht gut, nur bei ein Paar Sternen (wo auch die Luft etwas stark zitterte) sind sie [die Unterschiede] grösser, nämlich:
\begin{center}\begingroup\setlength{\tabcolsep}{4pt}
\begin{tabular}{@{}lrr|lrr@{}}
\(\eta\) Cephei&\(61^\circ8'30\overset{\prime\prime}{,}65\)&\(1{,}0\)&\(\omega'\) Cygni&\(48^\circ47'26\overset{\prime\prime}{,}06\)&\(1{,}8\)\\
&\(34{,}25\)&\(2{,}3\)&&\(25{,}17\)&\(2{,}3\)\\
&\(32{,}69\)&\(2{,}2\)&&\(28{,}74\)&\(0{,}9\)
\end{tabular}\par\endgroup\end{center}

\noindent\hspace*{\parindent}Hier nun meine Endresultate[\begingroup\renewcommand{\thefootnote}{*)}\footnotemark[\value{footnote}]\endgroup], denen ich noch das Datum, wofür das Mittel gilt, beigeschrieben habe, da die eigene Bewegung noch nicht abgezogen ist. Ich habe noch \(\alpha\) Lyrae und \(\alpha\) Cygni beigefügt:
\begin{center}\begingroup\fontsize{9.5}{12}\selectfont\setlength{\tabcolsep}{2pt}
\begin{tabular}{@{}lrrrclrlr@{}}
&&&&Gewicht&&&&\\
\(\eta\) Cephei&\(61^\circ\)&\(8'\)&\(32\overset{\prime\prime}{,}97\)&\(5{,}5\)&August&18.&B[\textsc{essel}:]&\(-1\overset{\prime\prime}{,}50\)\\
XX. 222&59&48&\(56{,}26\)&\(5{,}5\)&»&17.&&\(-0{,}84\)\\
\(\iota\) Draconis&59&36&\(0{,}19\)&\(4{,}5\)&»&20.&&—\\
47 »&59&10&\(15{,}18\)&\(5{,}3\)&Juli&21.&&\(+0{,}15\)\\
2 Ceph[ei] Hev.&58&22&\(4{,}48\)&\(4{,}9\)&August&9.&&\(+0{,}31\)\\
48 Draconis&57&34&\(44{,}79\)&\(6{,}1\)&Juli&21.&&\(-0{,}91\)\\
53 »&56&33&\(21{,}14\)&\(6{,}1\)&»&21.&&\(-0{,}62\)\\
33 Cygni&56&1&\(13{,}71\)&\(4{,}9\)&August&9.&&\(-0{,}86\)\\
49 Draconis&55&24&\(14{,}60\)&\(6{,}1\)&Juli&20.&&\(-0{,}78\)\\
46 »&55&21&\(38{,}26\)&\(6{,}1\)&»&21.&&\(-0{,}20\)
\end{tabular}\par\endgroup\end{center}
\clearpage\SourceLeaf{336}\HistoricalHead{331}{R}{BEOBACHTUNGEN AM REICHENBACHSCHEN MERIDIANKREISE.}\HistoricalFoot{}{42\textsuperscript{*}}
\begin{center}\begingroup\fontsize{9.5}{12}\selectfont\setlength{\tabcolsep}{2pt}
\begin{tabular}{@{}lrrrclrlr@{}}
XX. 391&53&49&\(48{,}37\)&\(4{,}5\)&August&20.&[\textsc{Bessel}:]&\(-0{,}74\)\\
51 Dracon[is]&53&7&\(27{,}35\)&\(6{,}1\)&Juli&21.&&\(-0{,}57\)\\
\(\kappa\) Cygni&53&2&\(25{,}99\)&\(6{,}1\)&»&21.&&\(-1{,}06\)\\
XXI. 32&52&49&\(56{,}04\)&\(4{,}5\)&August&19.&&\(-2{,}25\)\\
20 Cygni&52&32&\(6{,}67\)&\(5{,}0\)&»&9.&&\(-0{,}85\)\\
\(\iota\) »&51&21&\(2{,}15\)&\(5{,}7\)&»&4.&&\(-1{,}00\)\\
\(c\) » praec.&50&6&\(42{,}23\)&\(5{,}0\)&»&9.&&\(-0{,}54\)\\
\(c\) » sequ.&50&6&\(15{,}77\)&\(4{,}5\)&»&19.&&\(-1{,}10\)\\
\(\theta\) »&49&48&\(32{,}67\)&\(6{,}3\)&»&5.&&\(-1{,}19\)\\
\(\omega'\) »&48&47&\(26{,}13\)&\(5{,}0\)&»&12.&&\(-0{,}59\)\\
\(\alpha\) »&44&38&\(30{,}63\)&\(6{,}3\)&»&13.&&—\\
\(\alpha\) Lyrae&38&37&\(19{,}68\)&\(6{,}1\)&Juli&21.&&—\\\hline
\multicolumn{8}{c}{Mittel ohne Rücks[icht] auf Gew[icht]}&\(-0{,}80\)
\end{tabular}\par\endgroup\end{center}

\noindent\hspace*{\parindent}Von \(\alpha\) Cygni sind noch zwei alte Bestimmungen, welche geben
\begin{center}\(44^\circ38'30\overset{\prime\prime}{,}37\)\quad [Gewicht] \(2{,}0\)\quad Apr[il] 9.\end{center}
\noindent Auch ist früher \(\eta\) Cephei als Circumpolarstern in der ersten Periode beobachtet:
\begin{center}Decl[ination] \(61^\circ8'31\overset{\prime\prime}{,}92\)\quad Gewicht \(=4{,}8\)\quad Febr[uar] 29.\end{center}
\noindent Ebenso mehrere Ihrer andern Sterne, \(\delta\), \(\varepsilon\), \(\zeta\) Cephei und \(c\) Lacertae, deren Resultate aber heute der Kürze der Zeit wegen ich nicht ebenso zusammenstellen kann. So lasse ich auch erst künftig die Resultate für die \textsc{Maskelyne}schen Sterne vollständig folgen. . . . . . .

Ueber die Unterschiede unsrer Declinationen wage ich noch kein Urtheil zu äussern, wir müssen erst viel mehr Erfahrungen vergleichen. Wie glücklich würde man sich vor wenigen Jahren schon bei einer \textls[120]{solchen} Uebereinstimmung gehalten haben, zumal bei einem Instrument von so kleinen Dimensionen, was nicht repetirt und wo die Ablesung im Grunde noch der schwache Theil ist. . . . . . .

. . . . . . Damit Sie desto leichter prüfen können, ob vielleicht einer von uns bei Piazzi XXI. 32, der am meisten differirt, einen Reductionsfehler begangen hat, setze ich die Beob[achtungen] noch dafür einzeln her. . . . . . .
\clearpage\endgroup
